% GENERATED FILE. Edit canonical lessons, then run npm run generate:books.
% canonical-source-hash: fnv1a64:bf9af0452e255538
% canonical-lessons: MR-C09-ghene, MR-C09-vicharne, MR-C09-madat-karne, MR-C09-avadne

\chapter{Taking, Asking, Helping, Liking}
\label{ch:mr-doing-verbs}
\hlchaptermodality{12}

\begin{chapteropening}
\textbf{By the end of this chapter:} \emph{I can take, ask, help and say what I like and what I love in Marathi, build a fresh verb by putting \mr{करणे} behind a noun, and say why \mr{मला} \mr{मराठी} \mr{आवडते} has no room for me as its subject while \mr{मी} \mr{प्रेम} \mr{करतो} does.}

Say \mr{मला} \mr{मराठी} \mr{आवडते} and hear that you are not its subject, set it against \mr{मी} \mr{प्रेम} \mr{करतो} where you are, place both beside \mr{मला} \mr{मराठी} \mr{येते} and \mr{मला} \mr{मराठी} \mr{समजते} so one frame carries three meanings, and run the chapter's four verbs back with where each came from — Sanskrit \mr{ग्रह्} beside English grab, \mr{विचार} doing double duty as thought and question, Arabic \mr{मदद} through Persian, and \mr{आवडणे} with no Sanskrit ancestor to claim.
\end{chapteropening}

\section[gheṇe]{\mr{घेणे} (gheṇe) — \textquotedblleft{}to take,\textquotedblright{} which you have been saying since Chapter 7}

\label{lesson:MR-C09-ghene}



\begin{quote}
Say \textquotedblleft{}to write\textquotedblright{} and \textquotedblleft{}to read\textquotedblright{} — \emph{lihiṇe}, \emph{vāchṇe} — and what \emph{lihiṇe}'s root meant before writing existed. (\textbf{To scratch}.) This next verb has been in your mouth since the farewells, unrecognised.
\end{quote}

\subsection*{You'll want to know: \mr{घेणे}}
\begin{quote}
\textbf{\mr{घेणे}} — \emph{gheṇe} — \textbf{to take}
\end{quote}

\textbf{\mr{काळजी} \mr{घ्या}} — \emph{kāḷjī ghyā}, \textquotedblleft{}take care\textquotedblright{} — is this verb. \textbf{\mr{घ्या}} \emph{ghyā} is its respectful command, \textbf{\mr{घे}} \emph{ghe} the familiar one you would use with \textbf{\mr{तू}}. The plain present is \textbf{\mr{मी} \mr{घेतो}} \emph{mī gheto} / \textbf{\mr{घेते}} \emph{ghete}.

Note the \emph{-\mr{त}-} that appears in the present but not the infinitive: \emph{ghe-} on its own, \emph{ghe-t-o} when conjugated. Marathi's commonest verbs are the ones that bend.

\subsection*{The letters in this word}
\textbf{\mr{घ}} is \emph{gha} — a \emph{g} with a real puff of breath behind it, a separate letter from \textbf{\mr{ग}} \emph{ga}, exactly as \textbf{\mr{ख}} \emph{kha} is separate from \textbf{\mr{क}} \emph{ka}. Add the \emph{-e} stroke over the top and the retroflex \textbf{\mr{णे}}: \emph{ghe-ṇe}.

\begin{cousinweb}
\textbf{\mr{घेणे}} comes through Prakrit \emph{\textbf{gahaï}} from the Sanskrit root \textbf{\mr{ग्रह्}} \emph{grah-}, \textquotedblleft{}to seize, to grasp\textquotedblright{} — the root of \textbf{\mr{ग्रहण}} \emph{grahaṇ}, an eclipse, which is the sky being seized.

Its older Vedic shape was \textbf{\mr{ग्रभ्}} \emph{grabh-}, and that \textbf{bh} is the piece worth holding on to: Indo-European \textbf{*g\textsuperscript{h}reb\textsuperscript{h}-} kept its \emph{bh} in Vedic and turned it into a \emph{b} in Germanic, which is English \textbf{grab}. Marathi's \emph{gh-} is what four thousand years did to the same first consonant. \emph{Grabh} and \emph{grab} are the same word.
\end{cousinweb}

\subsection*{Your turn}
Say these aloud:

\begin{itemize}
\raggedright
  \item \textquotedblleft{}gheṇe\textquotedblright{} — to take — then \textquotedblleft{}mī gheto\textquotedblright{} or \textquotedblleft{}ghete\textquotedblright{}
  \item the farewell it was hiding in — \textquotedblleft{}kāḷjī ghyā\textquotedblright{}
  \item the Vedic form and its English cousin — \textquotedblleft{}grabh … grab\textquotedblright{}
\end{itemize}

\begin{tcolorbox}[breakable,colback=teal!4,colframe=teal!35!black,title={Before you move on}]
Which farewell has been using \textbf{\mr{घेणे}} all along, and what does it literally say? (\textbf{\mr{काळजी} \mr{घ्या}} — \textquotedblleft{}take care.\textquotedblright{}) Which English word is the same word as Vedic \emph{grabh}? (\textbf{Grab}.) What gender is \textbf{\mr{घेणे}} when it acts as a noun? (\textbf{Neuter}, like every \emph{-\mr{णे}} form.) Say \textquotedblleft{}to write\textquotedblright{} and \textquotedblleft{}to take\textquotedblright{} in one breath. (\emph{Lihiṇe, gheṇe}.)
\end{tcolorbox}
\section[vichārṇe]{\mr{विचारणे} (vichārṇe) — \textquotedblleft{}to ask,\textquotedblright{} built out of \textquotedblleft{}a thought\textquotedblright{}}

\label{lesson:MR-C09-vicharne}



\begin{quote}
Say \textquotedblleft{}to take\textquotedblright{} — \emph{gheṇe}. Say \textquotedblleft{}to think\textquotedblright{} — \emph{vichār karṇe} — and what its root meant. (\textbf{To turn, to range about}.) Hold that word. The next verb is it, with four letters fewer.
\end{quote}

\subsection*{You'll want to know: \mr{विचारणे}}
\begin{quote}
\textbf{\mr{विचारणे}} — \emph{vichārṇe} — \textbf{to ask}
\end{quote}

\textbf{\mr{मी} \mr{विचारतो}} — \emph{mī vichārto} — \textquotedblleft{}I ask,\textquotedblright{} \textbf{\mr{विचारते}} \emph{vichārte} if a woman speaks. And the same \textbf{\mr{च}} rule holds: say it \emph{vitsārto}, with the \emph{ts} Marathi puts before a long \emph{ā}.

\subsection*{The letters in this word}
There is nothing new. \textbf{\mr{विचारणे}} is \textbf{\mr{विचार}} — the noun from the last chapter, letter for letter — with \textbf{\mr{णे}} on the end. If you can write \textquotedblleft{}a thought,\textquotedblright{} you can write \textquotedblleft{}to ask.\textquotedblright{}

\begin{cousinweb}
Marathi builds both verbs on \textbf{\mr{विचार}}:

\noindent
\begin{tabularx}{\linewidth}{@{}>{\raggedright\arraybackslash}X>{\raggedright\arraybackslash}X@{}}
\toprule
\textbf{Marathi} & \textbf{what it means} \\
\midrule
\mr{विचार} \mr{करणे} & to think — \textquotedblleft{}to \textbf{do} a thought\textquotedblright{} \\
\mr{विचारणे} & to ask — the thought \textbf{made into a verb} \\
\bottomrule
\end{tabularx}

Both sit on \textbf{\mr{वि}-} plus \textbf{\mr{चर्}} \emph{car-}, \textquotedblleft{}to move, to range.\textquotedblright{} Thinking is ranging over something privately; asking is ranging over it out loud with somebody else. Marathi did not decide they were different enough to need different words.

This is Marathi's own arrangement, not a family inheritance. Hindi's verb for asking, \textbf{\mr{पूछना}} \emph{pūchnā}, comes from an unrelated root, and Hindi's \textbf{\mr{विचार}} stays a noun. Where a language draws its lines is a choice each language makes for itself.
\end{cousinweb}

\subsection*{Your turn}
Say these aloud:

\begin{itemize}
\raggedright
  \item the pair, slowly — \textquotedblleft{}vichār karṇe … vichārṇe\textquotedblright{}
  \item \textquotedblleft{}I ask\textquotedblright{} in your own gender — \textquotedblleft{}mī vichārto\textquotedblright{} or \textquotedblleft{}vichārte\textquotedblright{}
  \item both with the Marathi \emph{ts} — \textquotedblleft{}vitsār karto … vitsārto\textquotedblright{}
\end{itemize}

\begin{tcolorbox}[breakable,colback=teal!4,colframe=teal!35!black,title={Before you move on}]
What is the only difference between Marathi's \textquotedblleft{}to think\textquotedblright{} and its \textquotedblleft{}to ask\textquotedblright{}? (\textbf{\mr{करणे}} — thinking adds \textquotedblleft{}to do\textquotedblright{}; asking puts \emph{-\mr{णे}} straight onto \textbf{\mr{विचार}}.) What did the shared root mean? (\textbf{To move, to turn over}.) How is the \textbf{\mr{च}} said in both? (Nearer \textbf{ts} before the long \emph{ā}.) And \textquotedblleft{}to take\textquotedblright{}? (\emph{Gheṇe}.)
\end{tcolorbox}
\section[madat karṇe]{\mr{मदत} \mr{करणे} (madat karṇe) — \textquotedblleft{}to help,\textquotedblright{} a word that travelled}

\label{lesson:MR-C09-madat-karne}



\begin{quote}
Say \textquotedblleft{}to ask\textquotedblright{} — \emph{vichārṇe}. And say \textquotedblleft{}I know\textquotedblright{} — \emph{malā māhīt āhe} — then name where \textbf{\mr{माहीत}} was borrowed from. (\textbf{Arabic, through Persian}.) A second borrowed word from that same layer is coming.
\end{quote}

\subsection*{You'll want to know: \mr{मदत} \mr{करणे}}
\begin{quote}
\textbf{\mr{मदत} \mr{करणे}} — \emph{madat karṇe} — \textbf{to help}
\end{quote}

Built the way \emph{kām karṇe} and \emph{vichār karṇe} are built: a noun, \textbf{\mr{मदत}} \emph{madat} (\textquotedblleft{}help\textquotedblright{}), with \textbf{\mr{करणे}} behind it. \textbf{\mr{मी} \mr{मदत} \mr{करतो}} — \emph{mī madat karto} — \textquotedblleft{}I help,\textquotedblright{} \textbf{\mr{करते}} \emph{karte} for a woman.

That pattern is the most productive thing in the language: put \textbf{\mr{करणे}} behind almost any noun and Marathi will accept the verb.

\subsection*{The letters in this word}
\textbf{\mr{मदत}} is three bare consonants, each with its built-in \emph{a}: \textbf{\mr{म}} \emph{ma}, \textbf{\mr{द}} \emph{da}, \textbf{\mr{त}} \emph{ta} — \emph{ma-da-t}, with the last \emph{a} falling away in speech as it usually does at the end of a Marathi word.

\begin{cousinweb}
\textbf{\mr{मदत}} is \textbf{Arabic} \emph{madad}, borrowed through \textbf{Persian}. Its Arabic root means \textbf{to stretch out, to extend} — \emph{madad} is a thing extended: a hand, or reinforcements sent to an army. Help, named as reach.

It belongs to the layer of Persian and Arabic administration that settled into Marathi under the Deccan sultanates — the same layer that gave you \textbf{\mr{माहीत}} in \emph{malā māhīt āhe}, and \textbf{\mr{हरकत}} in \emph{kāhī harkat nāhī}. Three everyday words, one route.

One small tell worth keeping: Marathi writes \textbf{\mr{मदत}}, with a final \textbf{\mr{त}}. Hindi and Urdu write \textbf{\mr{मदद}}, with a \textbf{\mr{द}}. The same loan, landing differently in each language — which is what borrowing normally looks like.
\end{cousinweb}

\subsection*{Your turn}
Say these aloud:

\begin{itemize}
\raggedright
  \item \textquotedblleft{}madat karṇe\textquotedblright{} — then \textquotedblleft{}mī madat karto\textquotedblright{} or \textquotedblleft{}karte\textquotedblright{}
  \item the three \mr{करणे} verbs you have — \textquotedblleft{}kām karṇe, vichār karṇe, madat karṇe\textquotedblright{}
  \item what the Arabic root meant — \textquotedblleft{}to stretch out\textquotedblright{}
\end{itemize}

\begin{tcolorbox}[breakable,colback=teal!4,colframe=teal!35!black,title={Before you move on}]
What did \textbf{\mr{मदत}}'s Arabic root mean before it meant help? (\textbf{To stretch out, to extend}.) Which road did it travel to reach Marathi? (\textbf{Persian}.) Name another word you have from that same layer. (\textbf{\mr{माहीत}}, or \textbf{\mr{हरकत}} in \emph{kāhī harkat nāhī}.) And how does Marathi spell the end of the word, against Hindi and Urdu? (With \textbf{\mr{त}}, not \textbf{\mr{द}}.)
\end{tcolorbox}
\section[āvaḍṇe]{\mr{आवडणे} (āvaḍṇe) — \textquotedblleft{}to like,\textquotedblright{} which is not something you do}

\label{lesson:MR-C09-avadne}



\begin{quote}
Three verbs from this chapter: \emph{gheṇe}, \emph{vichārṇe}, \emph{madat karṇe}. Say them. Each let you be the one acting. The fourth will not.
\end{quote}

\subsection*{You'll want to know: \mr{आवडणे}}
\begin{quote}
\textbf{\mr{आवडणे}} — \emph{āvaḍṇe} — \textbf{to like}
\end{quote}

\textbf{\mr{मला} \mr{मराठी} \mr{आवडते}} — \emph{malā marāṭhī āvaḍte} — \textquotedblleft{}I like Marathi,\textquotedblright{} or word for word, \textquotedblleft{}\textbf{to me} Marathi pleases.\textquotedblright{}

\subsection*{The letters in this word}
\textbf{\mr{आ}} is the standalone long \emph{ā}; \textbf{\mr{व}} \emph{va}; \textbf{\mr{ड}} \emph{ḍa}, a \textbf{retroflex} \emph{d} made with the tongue curled back — new here, and not the \textbf{\mr{द}} \emph{da} of \emph{madat}; then \textbf{\mr{णे}}. \emph{Ā-va-ḍa-ṇe}.

\begin{grammarlens}[title={Grammar Lens: the third sentence built this way}]
\noindent
\begin{tabularx}{\linewidth}{@{}>{\raggedright\arraybackslash}X>{\raggedright\arraybackslash}X@{}}
\toprule
\textbf{Marathi} & \textbf{word for word} \\
\midrule
\mr{मला} \mr{मराठी} \mr{येते} & to-me Marathi comes — \textbf{I know it} \\
\mr{मला} \mr{मराठी} \mr{समजते} & to-me Marathi is-understood — \textbf{I understand it} \\
\mr{मला} \mr{मराठी} \mr{आवडते} & to-me Marathi pleases — \textbf{I like it} \\
\bottomrule
\end{tabularx}

One frame, three meanings. You are \textbf{\mr{मला}} in all three, never the subject, and the verb agrees with \textbf{\mr{मराठी}}. Change the thing liked and the ending follows it: a masculine thing takes \textbf{\mr{आवडतो}}, a neuter thing \textbf{\mr{आवडतं}}.
\end{grammarlens}

\begin{cousinweb}
\textbf{\mr{आवडणे}} is \textbf{Marathi's own} — straight down from Old Marathi, cousin to Konkani \emph{āvaḍ}, with \textbf{no Sanskrit ancestor securely established}, so this lesson claims none. Hindi and Urdu say \textbf{\mr{पसंद}} \emph{pasand}, borrowed from Persian.

Loving is a separate verb: \textbf{\mr{प्रेम} \mr{करणे}} \emph{prem karṇe}, and there you \textbf{are} the subject — \textbf{\mr{मी} \mr{प्रेम} \mr{करतो}}. \textbf{\mr{प्रेम}} is Sanskrit \emph{preman}, on \textbf{\mr{प्री}} \emph{prī-}, \textquotedblleft{}to be dear,\textquotedblright{} from Indo-European \textbf{*preyH-}. Germanic built a verb \textquotedblleft{}to love\textquotedblright{} on it, and from that English has \textbf{friend} — the one loved — and \textbf{free}.
\end{cousinweb}

\begin{culture}
Chapter 9's rule holds here: neither language is simply older. Marathi's \textbf{\mr{दोन}} took its final \emph{n} by copying \textbf{\mr{तीन}} — an innovation. Hindi's \textbf{\mr{पाँच}} kept a nasal Marathi dropped — a retention. Now a verb: Marathi kept its native word for liking, Hindi took a Persian one.
\end{culture}

\subsection*{Your turn}
Say these aloud:

\begin{itemize}
\raggedright
  \item \textquotedblleft{}malā marāṭhī āvaḍte\textquotedblright{} — hearing \textquotedblleft{}to me,\textquotedblright{} not \textquotedblleft{}I\textquotedblright{}
  \item the one you do — \textquotedblleft{}mī prem karto\textquotedblright{} or \textquotedblleft{}karte\textquotedblright{}
  \item this chapter's four — \textquotedblleft{}gheṇe, vichārṇe, madat karṇe, āvaḍṇe\textquotedblright{}
\end{itemize}

\begin{tcolorbox}[breakable,colback=teal!4,colframe=teal!35!black,title={Before you move on}]
Who is the subject of \emph{malā marāṭhī āvaḍte}? (\textbf{\mr{मराठी}} — not you.) Which of the two puts you in the subject's chair? (\textbf{\mr{प्रेम} \mr{करणे}}.) Which English words share \textbf{\mr{प्रेम}}'s root? (\textbf{Friend}, \textbf{free}.) Now the chapter's four sources — \emph{gheṇe}, \emph{vichārṇe}, \emph{madat karṇe}, \emph{āvaḍṇe}. (Sanskrit \textbf{\mr{ग्रह्}}, cousin of \emph{grab}; \textbf{\mr{विचार}}, the thinking word itself; \textbf{Arabic} \emph{madad} through Persian; and \textbf{\mr{आवड}}, native.)
\end{tcolorbox}
